% =============================================================================
% Patch: §5 Toy anomaly --- uniform-remainder gap closure
%   sec:toy-anomaly-and-transverse-obstruction
%
% Target file:   rh/rh_rebuild.tex (currently at master after the §5
%                migration patch rh/patches/sec-toy-anomaly-and-transverse-
%                obstruction.tex was applied).
% Action:        this patch is NOT to be merged into rh_rebuild.tex yet.
%                It is staged under rh/patches/ for review.
% Format:        purely additive.  Does not modify any existing
%                statement, proof, or remark in §5.  Inserts new content
%                only.
%
% Purpose:       discharge rem:wip-toy-anomaly-uniform-remainder by
%                stating the visibility lower bound on a
%                polynomial-weight retained subdomain
%                D_T^toy(D) ⊂ D_T, with an explicit density estimate.
%
% Honest scope:  the original wip asks for a uniform extension of
%                eq:toy-anomaly-visibility-A to all of D_T, including
%                d → 0, d → ∞, and d approaching a zero of F_∞.
%                Of these three regimes:
%                  - d = 0 is genuinely impossible:
%                    cor:omega-coalescence applied to the toy phase
%                    forces \widehat\Om_toy(0; m, u) = I_2 for every u,
%                    whence A_toy = 0 by lem:A-toy-baseline-vanishing;
%                    no u^2-lower bound can hold there.
%                  - d → ∞ and d near zeros of F_∞ are excluded by
%                    the compactness and avoidance hypotheses on D.
%                The polynomial-weight subdomain produced here is what
%                the §8.1 rigidity comparison
%                (subsec:two-point-comparison) consumes; the density
%                loss Q(T)/q(T) combines with the other polynomial
%                losses tracked in
%                sec:aggregation-and-zero-free-strip.
%
% Insertion points (relative to §5 of rh_rebuild.tex as currently
% committed):
%
%   (1) NEW SUBSECTION between subsec:toy-anomaly-functional and
%       subsec:toy-anomaly-visibility, containing
%       def:toy-retained-subdomain and lem:toy-retained-density.
%
%   (2) NEW COROLLARY at the end of subsec:toy-anomaly-visibility,
%       immediately after the proof of thm:toy-anomaly-visibility:
%       cor:toy-visibility-polynomial-weight.
%
%   (3) NEW REMARK appended at the end of §5, immediately after the
%       existing wip rem:wip-toy-anomaly-uniform-remainder:
%       rem:toy-anomaly-no-uniform-extension.  This is a closed
%       (non-wip) remark that documents which sub-claims of the wip
%       are genuinely impossible (d = 0 by coalescence) and which are
%       handled by the polynomial-weight scope established here.
%       The wip itself remains in place; once a referee or coordinator
%       agrees that the new content discharges it, the wip can be
%       removed in a separate edit.
%
% Closes (when merged + wip removed):
%   rem:wip-toy-anomaly-uniform-remainder.
%
% Status (when merged):
%   \statusmig      partial -> migrated  (after wip removal)
%   \statusreferee  in-progress -> in-progress (re-review needed)
%   \statussympy    [..] verified -> [..] verified (extended via
%                     rh/sympy/sec-toy-anomaly-uniform-remainder-closure.py)
%   \statussim      [..] verified -> [..] verified (extended via
%                     rh/simulation/sec-toy-anomaly-uniform-remainder-closure.py)
%
% Verification (every claim has a failing test in the cited script):
%   - rh/sympy/sec-toy-anomaly-uniform-remainder-closure.py
%       verify_toy_retained_density_inequality_symbolic,
%       verify_density_constant_positivity_on_sample_D,
%       verify_height_hypothesis_consistency.
%   - rh/simulation/sec-toy-anomaly-uniform-remainder-closure.py
%       check_density_at_heights,
%       check_density_scaling_in_Q,
%       check_height_hypothesis_holds_at_high_T,
%       check_avoidance_set_for_F_inf_zeros.
% =============================================================================

% =============================================================================
% (1) Insert this subsection BETWEEN subsec:toy-anomaly-functional
%     (which currently ends with the proof of lem:toy-anomaly-leading
%     in rh_rebuild.tex) AND subsec:toy-anomaly-visibility (which
%     currently begins with thm:toy-anomaly-visibility).
% =============================================================================

\subsection{Toy retained subdomain and polynomial-weight density}
\label{subsec:toy-retained-subdomain}

\begin{definition}[Toy retained subdomain]
\label{def:toy-retained-subdomain}
Fix a compact interval \(D=[d_-,d_+]\subset(0,\infty)\) avoiding the
isolated zeros of \(F_{\infty}\) of
Lemma~\ref{lem:toy-anomaly-leading}.  The toy retained subdomain at
height \(T\) is
\begin{equation}
\label{eq:toy-retained-subdomain}
        \mathcal D_T^\toy(D)
        \;:=\;\bigl\{(m,s)\in\mathcal D_T^\times :
            q(m)\,s\in D\bigr\}.
\end{equation}
\end{definition}

\begin{lemma}[Polynomial-weight density of the toy retained subdomain]
\label{lem:toy-retained-density}
For every compact \(D=[d_-,d_+]\subset(0,\infty)\), at every height
\(T\) with
\begin{equation}
\label{eq:toy-density-height-hypothesis}
        \frac{d_+}{q(T)}\;\le\;\frac{|I_T|}{4},
\end{equation}
the toy retained subdomain satisfies
\begin{equation}
\label{eq:toy-retained-density}
    \frac{|\mathcal D_T^\toy(D)|}{|\mathcal D_T|}
    \;\ge\;
    \frac{(d_+-d_-)\,Q(T)}{q(T)},
\end{equation}
where \(|\cdot|\) denotes two-dimensional Lebesgue measure on the
\((m,s)\) plane.  The right-hand side is positive and
polynomial-in-\(Q\) under any retained \(Q\)-scaling with
\(\log T\le Q(T)^{c}\) for some absolute \(c\ge 1\).  The height
hypothesis \eqref{eq:toy-density-height-hypothesis} is equivalent to
\(Q(T)/q(T)\le 1/(4 d_+)\) under the chart convention
\(|I_T|\le 1/Q(T)\) used in
Section~\ref{sec:local-kernel-and-jet-normalization}; it is satisfied
for all sufficiently large \(T\) whenever \(Q(T)=o(q(T))\), which is
the slowly-growing-\(Q\) regime of
Section~\ref{subsec:notation-and-conventions}.
\end{lemma}

\begin{proof}
Recall that
\[
    \mathcal D_T
    \;=\;\bigl\{(m,s)\in\mathbb R^{\,2}: t_-=m-s/2\in I_T,\;
        t_+=m+s/2\in I_T\bigr\},
\]
which is the rotated square
\(\{(m,s):|m-T|+|s|/2\le |I_T|/2\}\) of area \(|\mathcal D_T|=|I_T|^{2}\).

Restrict to the slab \(|m-T|\le|I_T|/4\).  For such \(m\), the
\(\mathcal D_T\) constraint reduces to \(|s|\le|I_T|/2-|m-T|\ge|I_T|/4\).
By Lemma~\ref{lem:theta-derivative-asymptotics},
\(q(m)\le q(T)\) on \(I_T\), so the toy constraint
\(q(m)\,s\in D\) reads
\(s\in[d_-/q(m),d_+/q(m)]\), an interval of length
\((d_+-d_-)/q(m)\ge(d_+-d_-)/q(T)\).  Under
\eqref{eq:toy-density-height-hypothesis},
\(d_+/q(m)\le d_+/q(T)\cdot 1\le|I_T|/4\), so this \(s\)-interval is
contained in \([0,|I_T|/4]\) and hence in the allowed
\(\mathcal D_T\)-range \([0,|I_T|/2-|m-T|]\) for every
\(|m-T|\le|I_T|/4\).  By symmetry the negative-\(s\) interval
\([-d_+/q(m),-d_-/q(m)]\) contributes equally.  Therefore
\[
    |\mathcal D_T^\toy(D)|
    \;\ge\;2\int_{\{|m-T|\le|I_T|/4\}}\frac{d_+-d_-}{q(T)}\,dm
    \;=\;\frac{(d_+-d_-)\,|I_T|}{q(T)}.
\]
Comparing with \(|\mathcal D_T|=|I_T|^{2}\),
\[
    \frac{|\mathcal D_T^\toy(D)|}{|\mathcal D_T|}
    \;\ge\;\frac{d_+-d_-}{q(T)\,|I_T|}.
\]
The convention \(|I_T|\le 1/Q(T)\) of
Section~\ref{sec:local-kernel-and-jet-normalization}
gives \eqref{eq:toy-retained-density}.  For the equivalence stated
after the lemma, multiply the hypothesis
\eqref{eq:toy-density-height-hypothesis} on both sides by \(4/|I_T|\)
and use \(1/|I_T|\ge Q(T)\) to obtain \(4 d_+ Q(T)/q(T)\le 1\).
\end{proof}

% =============================================================================
% (2) Insert this corollary at the END of subsec:toy-anomaly-visibility,
%     immediately after the proof of thm:toy-anomaly-visibility.
% =============================================================================

\begin{corollary}[Visibility on the polynomial-weight subdomain]
\label{cor:toy-visibility-polynomial-weight}
Fix \(D\) and the constants \(Q_\toy(D)\), \(u_\toy(D)\), \(c_\toy(D)\)
of Theorem~\ref{thm:toy-anomaly-visibility}.  For every
\(0<|u|\le u_\toy(D)\), every \(q_0\ge Q_\toy(D)\), and every
\((m,s)\in\mathcal D_T^\toy(D)\) of
Definition~\ref{def:toy-retained-subdomain},
\begin{equation}
\label{eq:toy-anomaly-visibility-on-subdomain}
        \bigl|\mathcal A_\toy\bigl(\widehat\Om_\toy(s;m,u)\bigr)\bigr|
        \;\ge\;\tfrac12\,c_\toy(D)\,u^{2}.
\end{equation}
By Lemma~\ref{lem:toy-retained-density},
\(\mathcal D_T^\toy(D)\) carries a relative Lebesgue density at
least \((d_+-d_-)\,Q(T)/q(T)\) inside \(\mathcal D_T\).  Hence the
visibility lower bound is consumed by
Section~\ref{sec:projective-rigidity} on a polynomial-weight subset
of midpoints and separations, with the same exponent \(A_\toy=0\) and
the same constant \(c_\toy(D)/2\), and a polynomial-in-\(Q\)
density factor that combines with the other polynomial losses tracked
in Section~\ref{sec:aggregation-and-zero-free-strip}.
\end{corollary}

\begin{proof}
For \((m,s)\in\mathcal D_T^\toy(D)\) one has \(q(m)\,s\in D\) by
\eqref{eq:toy-retained-subdomain}; with \(q_0=q(m)\),
\(d=q_0\,s\in D\).  Theorem~\ref{thm:toy-anomaly-visibility} then
gives \eqref{eq:toy-anomaly-visibility-on-subdomain}.  The density
estimate is Lemma~\ref{lem:toy-retained-density}.
\end{proof}

% =============================================================================
% (3) Append this remark at the END of §5, immediately after the
%     existing wip rem:wip-toy-anomaly-uniform-remainder.  The remark
%     documents why the wip's "uniform extension to all of D_T" is
%     genuinely impossible at d = 0 and how the polynomial-weight
%     scope above subsumes the other regimes mentioned in the wip.
% =============================================================================

\begin{remark}[No uniform extension; polynomial-weight scope]
\label{rem:toy-anomaly-no-uniform-extension}
The visibility lower bound \eqref{eq:toy-anomaly-visibility-A} cannot
be extended to all of \(\mathcal D_T\).  At \(s=0\),
Corollary~\ref{cor:omega-coalescence} applied to the toy phase gives
\(\widehat\Om_\toy(0;m,u)=I_{2}\) for every \(u\), whence
\(\mathcal A_\toy(\widehat\Om_\toy(0;m,u))=0\) by
Lemma~\ref{lem:A-toy-baseline-vanishing}; no \(u^{\,2}\)-lower bound
can hold at the coalescence locus \(d=0\).  The regimes \(d\to\infty\)
and \(d\) approaching an isolated zero of \(F_{\infty}\) are excluded
by the compactness and avoidance hypotheses on \(D\) of
Definition~\ref{def:toy-retained-subdomain}.  The polynomial-weight
subdomain \(\mathcal D_T^\toy(D)\) of
Lemma~\ref{lem:toy-retained-density} carries
\eqref{eq:toy-anomaly-visibility-A} via
Corollary~\ref{cor:toy-visibility-polynomial-weight}; the density
factor \(Q(T)/q(T)\) combines with the other polynomial losses tracked
in Section~\ref{sec:aggregation-and-zero-free-strip}.  This
polynomial-weight statement is what the comparison theorems of
Section~\ref{sec:projective-rigidity} consume.
\end{remark}

% =============================================================================
% End of patch.
% =============================================================================
